%======================================================================
% Plantilla de CV — Juan García López
%
% Marcadores que rellena el pipeline (NO los borres si editas el diseño).
% Se escriben entre corchetes en el cuerpo del documento; aquí se nombran
% sin corchetes para que la sustitución no toque estos comentarios:
%   HEADLINE         titular bajo el nombre (p. ej. INGENIERO DE PROCESOS)
%   PROFILE_SUMMARY  párrafo "Mi perfil" adaptado a la oferta
%   INERCO_BULLETS   viñetas \item de la experiencia en INERCO
%   COMPANY_NAME y JOB_TITLE  se usan en los metadatos del PDF
%
% La foto debe llamarse foto.jpg y estar en plantillas/assets/.
%======================================================================
\documentclass[10pt,a4paper]{article}
\usepackage[utf8]{inputenc}
\usepackage[T1]{fontenc}
\usepackage{lmodern}
\usepackage[a4paper,margin=0cm]{geometry}
\usepackage{graphicx}
\usepackage{xcolor}
\usepackage{tikz}
\usepackage{enumitem}
\usepackage{microtype}
\usepackage[hidelinks]{hyperref}
\hypersetup{pdftitle={CV Juan García López - [JOB_TITLE] - [COMPANY_NAME]}}

\definecolor{azulcv}{HTML}{1B3155}
\definecolor{grisfondo}{HTML}{E9EBF0}
\definecolor{grisoscuro}{HTML}{555C66}

\renewcommand{\familydefault}{\sfdefault}
\pagestyle{empty}
\setlength{\parindent}{0pt}

% ---------- estilos ----------
\newcommand{\tituloseccion}[1]{%
  {\color{azulcv}\Large\bfseries\textls[60]{#1}}\\[2pt]
  {\color{azulcv}\rule{2.2cm}{1.6pt}}\\[10pt]}
\newcommand{\titulolateral}[1]{%
  {\color{azulcv}\large\bfseries\textls[60]{#1}}\\[6pt]}
\newcommand{\entrada}[2]{%
  {\bfseries\color{azulcv}#1}\\[1pt]
  {\small\color{grisoscuro}#2}\\[4pt]}
\setlist[itemize]{leftmargin=12pt,itemsep=2pt,parsep=0pt,topsep=2pt,label=\textbullet}

\begin{document}

% ============================ PÁGINA 1 ============================
\begin{tikzpicture}[remember picture,overlay]
  % Banner superior azul
  \fill[azulcv] (current page.north west) rectangle ([yshift=-4.6cm]current page.north east);
  % Columna lateral gris
  \fill[grisfondo] ([yshift=-4.6cm]current page.north west)
    rectangle ([xshift=7.0cm]current page.south west);

  % Foto con marco blanco (solapada sobre el banner)
  \node[anchor=north west,inner sep=0]
      at ([xshift=1.15cm,yshift=-1.05cm]current page.north west)
    {\setlength{\fboxsep}{3pt}\colorbox{white}{\includegraphics[width=3.2cm]{foto.jpg}}};

  % Nombre y titular adaptado a la oferta
  \node[anchor=north east,inner sep=0,text width=12.2cm,align=flush right]
      at ([xshift=-1.3cm,yshift=-1.25cm]current page.north east) {%
    {\color{white}\fontsize{25}{30}\selectfont\bfseries JOSÉ MANUEL\\[2pt] CORONEL MORENO}\\[14pt]
    {\color{white}\normalsize\bfseries\textls[220]{[HEADLINE]}}%
  };

  % ------------------------ Columna lateral ------------------------
  \node[anchor=north west,inner sep=0,text width=4.9cm]
      at ([xshift=1.05cm,yshift=-6.4cm]current page.north west) {%
    \titulolateral{CONTACTO}
    \small
    \textbf{Teléfono:} 672 871 128\\[2pt]
    \textbf{Email:} \href{mailto:juangarcialopez@gmail.com}{juangarcialopez@gmail.com}\\[2pt]
    \textbf{Ubicación:} Huelva, España\\[16pt]
    \titulolateral{INFORMACIÓN}
    \small
    \begin{itemize}
      \item Carnet de conducir
      \item Disponibilidad horaria
      \item Edad: 22 años
    \end{itemize}
    \vspace{12pt}
    \titulolateral{HERRAMIENTAS}
    \small
    Excel \,·\, Python \,·\, MATLAB\\[2pt]
    PVsyst \,·\, SAM \,·\, AutoCAD\\[2pt]
    Symmetry\\[16pt]
    \titulolateral{IDIOMAS}
    \small
    \textbf{Español:} nativo\\[2pt]
    \textbf{Inglés:} B2 Cambridge%
  };

  % ----------------------- Columna principal -----------------------
  \node[anchor=north west,inner sep=0,text width=12.1cm]
      at ([xshift=8.0cm,yshift=-5.5cm]current page.north west) {%
    \tituloseccion{MI PERFIL}
    [PROFILE_SUMMARY]\\[16pt]
    \tituloseccion{FORMACIÓN}
    \entrada{INGENIERÍA INDUSTRIAL (2022 -- 2026)}
            {E.T.S.I. -- Universidad de Sevilla \,·\, Nota media: 8.0 \,·\, Nota TFG: 10}
    \entrada{CURSO EN INDUSTRIA QUÍMICA (2026)}
            {Grupo Aspasia -- Códice Cantabria}
    \entrada{CURSO EN ANÁLISIS Y PREVENCIÓN DE RIESGOS LABORALES (2026)}
            {E.T.S.I. -- Universidad de Sevilla}
    \entrada{GRADO DE ADMINISTRACIÓN (2016 -- 2022)}
            {IES la Orden -- Huelva}
    \vspace{12pt}
    \tituloseccion{EXPERIENCIA}
    \entrada{BECARIO DE INGENIERÍA -- MOEVE}{Febrero 2026 -- Agosto 2026}
    \begin{itemize}
[INERCO_BULLETS]
    \end{itemize}
    \vspace{8pt}
    \entrada{PROFESOR DE QUIMICA}{2023 -- 2025}
    \begin{itemize}
      \item Enseñanza de química
            a alumnos de entre 6 y 12 años.
    \end{itemize}
    \vspace{8pt}
    \entrada{PROFESOR DE CLASES PARTICULARES}{Julio y agosto de 2023}
    \begin{itemize}
      \item Impartí asignaturas como Física, Matemáticas o Inglés a alumnos
            de ESO y Bachillerato.
    \end{itemize}%
  };
\end{tikzpicture}
\null

% ============================ PÁGINA 2 ============================
\newpage
\begin{tikzpicture}[remember picture,overlay]
  % Banda superior azul
  \fill[azulcv] (current page.north west) rectangle ([yshift=-2.6cm]current page.north east);
  \node[anchor=west,inner sep=0] at ([xshift=1.2cm,yshift=-1.3cm]current page.north west)
    {\color{white}\Large\bfseries JUAN GARCÍA LÓPEZ};
  \node[anchor=east,inner sep=0] at ([xshift=-1.2cm,yshift=-1.3cm]current page.north east)
    {\color{white}\normalsize\bfseries\textls[160]{[HEADLINE]}};

  % Contenido a página completa
  \node[anchor=north west,inner sep=0,text width=18.6cm]
      at ([xshift=1.2cm,yshift=-3.8cm]current page.north west) {%
    \tituloseccion{PROYECTOS ACADÉMICOS}
    \entrada{INSTALACIÓN FOTOVOLTAICA DE AUTOCONSUMO EN UN INSTITUTO}
            {Diseño completo y simulación en PVsyst}
    \begin{itemize}
      \item Dimensioné la instalación según el consumo real del centro y la radiación
            de la localización.
      \item Estudié el retorno económico de la instalación con baterías y sin ellas.
      \item Seleccioné los componentes: panel fotovoltaico, inversor y batería.
      \item Simulé la instalación completa con el programa PVsyst.
    \end{itemize}
    \vspace{10pt}
    \entrada{DISEÑO DE CENTRAL TERMOSOLAR CILINDROPARABÓLICA}
            {Modelado con SAM (System Advisor Model)}
    \begin{itemize}
      \item Diseñé la central cilindroparabólica situada en San José del Valle con SAM.
      \item Estudié la viabilidad del HTF tanto con aceite térmico como con sales fundidas.
      \item Optimicé las horas de almacenamiento y el múltiplo solar.
    \end{itemize}
    \vspace{10pt}
    \entrada{VIABILIDAD TÉCNICA DE UN PARQUE EÓLICO OFFSHORE EN EL ESTRECHO DE GIBRALTAR}
            {Estudio de emplazamiento y tecnología}
    \begin{itemize}
      \item Seleccioné el emplazamiento atendiendo a profundidades, zonas restringidas
            y velocidades de viento.
      \item Elegí la turbina más apropiada para el parque y el tipo de cimentación.
      \item Localicé el punto de conexión en tierra.
    \end{itemize}
    \vspace{10pt}
    \tituloseccion{RECONOCIMIENTOS}
    \begin{itemize}
      \item Matrícula de Honor en Bachillerato.
    \end{itemize}%
  };
\end{tikzpicture}
\null
\end{document}
